\scanpage{29}
% Top.
Donc on trouve un foncteur \deleted{\unsure{contravariant}}
\[
\mathcal V\longrightarrow\text{Algèbres bivariantes jaugées}.
\]
3) Dans $\mathcal V^\circ$ \deleted{\unsure{il y a}} il y a objet initial: $k$.

Si $A,B\in\mathcal V^\circ$, il existe un coproduit dans $\mathcal V^\circ$ (\unsure{produit} de $\mathcal V$) savoir le produit tensoriel gauche $A\otimes B$. Soit
\[
f^*:A\mathbin{\text{---}}A',\quad g^*:B\mathbin{\text{---}}B',\quad\text{alors}\quad
f^*\otimes g^*:A\otimes B\mathbin{\text{---}}A'\otimes B'
\]
défini par
\[
f^*\otimes g^*(a\otimes b)=f^*(a)\otimes g^*(b)
\]
est l'hom correspondant des coproduits. \deleted{\unsure{On le note}} \inserted{On convient de} mettre sur $A\otimes B$ l'augmentation définie par
\[
\unsure{k|\kappa}_{A\otimes B}=\unsure{k|\kappa}_A\otimes\unsure{k|\kappa}_B
\]
(produit tensoriel \uline{gauche} d'homs homogènes d'\unsure{espaces} gradués).
\margin{i.e.\quad $\unsure{k|\kappa}_{A\otimes B}(x\otimes y)=\unsure{k|\kappa}_A(x)\unsure{k|\kappa}_B(\unsure{x|y})(-1)^{Bx}$}
\margin{produit tensoriel \uline{gauche} d'\unsure{homom.} \uline{homogènes} des \deleted{groupes}\inserted{modules} gr.dués}
\deleted{\unsure{On démontre que}}\quad\inserted{On a alors}
\[
\varphi(f^*\otimes g^*)=\varphi(f^*)\otimes\varphi(g^*)
\]
\margin{\unsure{2|Z}}
i.e.
\[
(*)\qquad(f\times g)_*(a'\otimes b')=f_*(a')\otimes g_*(b')(-1)^{ga'}.
\]
\figure{\begin{tikzpicture}\draw (0,2)--(1,2) .. controls (1.5,2) and (.8,1.5) .. (1.5,1.5)--(5,1.7) .. controls (5.8,1.7) and (5.2,1.2) .. (6,1.3)--(9,1.5);\draw[->] (1.8,.1)..controls (.5,.8) and (1,1.4)..(1.5,1.5);\draw (1.9,-.2)--(1.5,-.2)--(1.5,1.3);\end{tikzpicture}}

% Retraction and replacement signs.
\uline{Mais c'est faux} si on ne suppose pas $f$ et $g$ de degrés pairs.

Les deux mb. de la formule $(*)$ diffèrent par un facteur
\[
(-1)^{ag+bf+Bf}.
\]
Si dans la définition de $\unsure{k|\kappa}_{A\otimes B}$ on n'avait pas mis le \uline{signe} $(-1)^{Bx}$, le signe qui viendrait avec ces \unsure{not.}
\[
(-1)^{a'g+bf}.
\]
Si \inserted{en plus} \unsure{on se dispensait du $\otimes$ du produit|on disposait du $\otimes$ du produit} $f_*\otimes g_*$ au sens \deleted{gauche}\inserted{tordu}, et qu'on \deleted{\unsure{pour}} enlève le \uline{signe} $(-1)^{ga'}$; on trouverait donc un signe
\[
(-1)^{bf}=(-1)^{B'f}.
\]
Si on prenait $\unsure{k|\kappa}$ tordu, $\otimes$ ordinaire, le \uline{signe} serait
\[
(-1)^{bf+(AB+A'B')}.
\]
C'est donc le signe où on prend tout ordinaire
\[
(\unsure{k|\kappa}_{A\otimes B}\inserted{{}=\unsure{k|\kappa}_A\otimes\unsure{k|\kappa}_B}\text{ et }f_*\otimes g_*)
\]
qui est le plus simple. Donc on trouve alors
\[
\boxed{(f\times g)_*=(f_*\overset{\mathrm{ord}}\otimes g_*)
(\mathrm{id}_A\otimes\unsure{[\deleted{(-1)^B}(-1)_{B'}]^{\deg f}})}
\]
\margin{\deleted{\unsure{NB}} Il n'y a aucun signe dès que \textcircled{$f$} est pair.}
\mspage{28}
